% !TeX spellcheck = en_US
\documentclass[french]{yLectureNote}

\title{Algèbre linéaire 2}
\subtitle{Langage mathématique}
\author{Paulhenry Saux}
\date{\today}
\yLanguage{Français}

\professor{C.Dartyge}

\usepackage{graphicx}%----pour mettre des images
\usepackage[utf8]{inputenc}%---encodage
\usepackage{geometry}%---pour modifier les tailles et mettre a4paper
%\usepackage{awesomebox}%---pour les boites d'exercices, de pbq et de croquis ---d\'esactiv\'e pour les TP de PC
\usepackage{tikz}%---pour deiffner + d\'ependance de chemfig
\usepackage{tkz-tab}
\usepackage{chemfig}%---pour deiffner formules chimiques
\usepackage{chemformula}%---pour les formules chimiques en \'equation : \ch{...}
\usepackage{tabularx}%---pour dimensionner automatiquement les tableaux avec variable X
\usepackage{awesomebox}%---Pour les boites info, danger et autres
\usepackage{menukeys}%---Pour deiffner les touches de Calculatrice
\usepackage{fancyhdr}%---pour les en-t\^ete personnalis\'ees
\usepackage{blindtext}%---pour les liens
\usepackage{hyperref}%---pour les liens (\`a mettre en dernier)
\usepackage{caption}%---pour la francisation de la l\'egende table vers Tableau
\usepackage{pifont}
\usepackage{array}%---pour les tableaux
\usepackage{lipsum}
\usepackage{yFlatTable}
\usepackage{multicol}
\newcommand{\Lim}[1]{\lim\limits_{\substack{#1}}\:}
\renewcommand{\vec}{\overrightarrow}
\newcommand{\bz}{\overline{z}}
\DeclareMathOperator{\card}{card}
\renewcommand{\vec}{\overrightarrow}
\newcommand{\N}[0]{\mathbb{N}}
\newcommand{\R}[0]{\mathbb{R}}
\newcommand{\C}[0]{\mathbb{C}}
\newcommand{\dd}[0]{\mathrm{d}}
\begin{document}
\chapter{Déterminant}
\section{Une forme multilinéaire}
\begin{theorem}[Application linéaire du déterminant pour les colonnes]
Le déterminant est une application linéaire pour chaque colonne d'un ematrice :
\begin{itemize}
 \item \(\det||c_1, \dots, \lambda c_k, \dots, c_n|| = \lambda \det||c_1, \dots, c_k, \dots, c_n||\)
 \item \(\det||c_1, \dots, c_j + c_k, \dots, c_n|| =  \det||c_1, \dots, c_k, \dots, c_n|| + \det||c_1, \dots, c_j, \dots, c_n||\)
\end{itemize}
\end{theorem}
\begin{proposition}[Échange de colonne]
Si on échange de 2 colonnes, le déterminant change de signe
\end{proposition}
\begin{theorem}[Base et déterminant]
 Soit \(A = ||c_1, \dots, c_n||\). Ces vecteurs forment un ebase de $K^n$ ssi \(\det(A)\neq 0\).
\end{theorem}
\section{Permutations, transposition}
\begin{definition}[Transposiiton]
C'est une permutation qui échange juste 2 éléments différentes. les autres valaurs restent inchangées.
\end{definition}
\begin{proposition}[Décomposition de permutation]
Toute permutation peut se décomposer en un produit de transpositions.
\end{proposition}
\section{Déterminant de la transposée}
\begin{theorem}[Déterminant de la transposée]
 \(\det(^tA) = \det(A)\)
\end{theorem}
\section{Calcul de déterminants}
\begin{definition}[Cofacteur]
\(cof(a_{ij}) = (-1)^{i+j}\det(A_{ij})\)
\end{definition}
\begin{proposition}[Combinaison linéaire et déterminant]
Le déterminant ne change pas si à une ligne (respect. à une
colonne) on ajoute une combinaison linéaire des autres lignes (respect. des autres
colonnes).
\end{proposition}
\section{Déterminant de matrice}
\begin{theorem}[Produit de déterminant]
 \(\det(AB) = \det(A)\det(B)\)
\end{theorem}
\begin{theorem}
A est  est inversible si et seulement si \(\det A \neq 0\)
Si A est inversible : \(\det(A^{-1}) = \frac{1}{\det(A)})\)
\end{theorem}
\section{Calcul de l'inverse d'une matrice}
\begin{theorem}[Inverse d'une matrice 2]
 \(A^{-1} = (ad-bc)^{-1}(d -b // -c a)\)
\end{theorem}
\section{Application aux familles}
\begin{theorem}[Base et déterminant]
 Les vecteur colonnes forment une base si et seulement si \(\det \neq 0\)
\end{theorem}
\begin{definition}[Mineur]
On appelle mineur d’ordre r d’une matrice A le déterminant d’une matrice d’ordre
r extraite de A obtenue en choisissant r lignes et r colonnes. Ces lignes et colonnes peuvent ne pas \^etre consécutives.
\end{definition}
\begin{theorem}[Famille libre et mineur]
Soient r vecteurs d’un espace vectoriel E de dimension n (r \(\leq\) n) et
A la matrice dont les colonnes sont les composantes des vecteurs dans une base quelconque de E.
La famille est libre si et seulement si on peut extraire de A un mineur
d’ordre r non nul.
\end{theorem}
\begin{theorem}[Rang et mineur]
 Le rang d'une matrice A est l'ordre max des mineuys non nuls extraits de A, i.e. \(rg(A) = r \iff\) il existe un mineur d'ordre r non nul et tous les mineur d'ordre supéréiru sont nuls.
\end{theorem}

\end{document}


